\begin{titlepage}
\thispagestyle{empty}

\begin{center}

% ───────────────── HEADER ─────────────────
\begin{minipage}[t]{0.32\textwidth}
\raggedright
\small
République Tunisienne\\
Ministère de l'Enseignement\\
Supérieur et de la Recherche\\
Scientifique

\vspace{0.1cm}
\rule{2.5cm}{0.4pt}

\vspace{0.1cm}

Université de Gafsa\\
Institut Supérieur des Sciences\\
Appliquées et de Technologie\\
de Gafsa
\end{minipage}
\hfill
\begin{minipage}[t]{0.28\textwidth}
\centering
\includegraphics[width=2.3cm]{figures/40740c82-2e26-4178-bf35-c7f001b6e764.jpeg}\\[0.15cm]
\includegraphics[width=1.8cm]{figures/43151050-7e39-4248-a071-b25d81e3094d.jpeg}
\end{minipage}
\hfill
\begin{minipage}[t]{0.32\textwidth}
\raggedleft
\small
Cycle de Formation en Mastère\\
professionnel\\
dans la discipline\\
............................

\vspace{0.1cm}
\rule{2.5cm}{0.4pt}

\vspace{0.1cm}

\textbf{Mémoire de MASTÈRE}\\
D'ordre : 2025 -- ESC
\end{minipage}

\vspace{0.2cm}

\color{blue}
\rule{\textwidth}{1pt}
\color{black}

\vspace{0.2cm}

% ───────────────── TITRE ─────────────────
{\fontsize{22}{24}\selectfont \textbf{MEMOIRE}}

\vspace{0.3cm}

\small
\textit{Présenté à}\\[0.15cm]
L'Institut Supérieur des Sciences Appliquées et de Technologie de Gafsa

\vspace{0.3cm}

Département : Informatique

\vspace{0.3cm}

\textit{En vue de l'obtention du diplôme de}

\vspace{0.3cm}

{\Large \textbf{MASTÈRE PROFESSIONNEL}}
\vspace{0.3cm}
\begin{flushleft}
\hspace{2cm}\textbf{Mention :} \dotfill
\end{flushleft}

\vspace{-0.1cm}

\begin{flushleft}
\hspace{2cm}\textbf{Spécialité :} Sciences des données
\end{flushleft}

\vspace{0.2cm}

\textit{Par :}\\[0.1cm]
{\large \textbf{Marwa ABADA}}

\vspace{0.2cm}

\rule{7cm}{0.8pt}

\vspace{0.2cm}

{\normalsize \bfseries
Analyse de données médicales massives pour la prédiction\\[0.15cm]
des maladies cardiovasculaires à l'aide d'Apache Spark\\[0.15cm]
et des réseaux de neurones profonds
}

\vspace{0.2cm}

\rule{7cm}{0.8pt}

\vspace{0.2cm}

\small
Soutenu le .../.../2025 devant le jury composé de :

\vspace{0.3cm}

\begin{tabular}{p{7.5cm}p{3.5cm}}
M. ........................................ & Président \\
M. ........................................ & Rapporteur \\
M. ........................................ & Encadrant \\
\end{tabular}

\vfill

{\large \textbf{A.U : 2024 -- 2025}}

\end{center}
\end{titlepage}

% ─────────────────────────────────────────────────────────────────────────────
% DÉDICACE
% ─────────────────────────────────────────────────────────────────────────────
\thispagestyle{empty}
\null\vfill
\begin{center}
  {\color{black}\rule{0.6\textwidth}{1pt}}\\[1cm]
  {\Large\itshape\bfseries Dédicace}\\[1cm]
  {\color{black}\rule{0.6\textwidth}{0.4pt}}
\end{center}
\vspace{1cm}

\begin{center}
\begin{minipage}{0.75\textwidth}
\setlength{\parindent}{0pt}
\setlength{\parskip}{0.8em}
\itshape

À mes très chers parents,\\
pour leur amour inconditionnel, leurs sacrifices constants\\
et leur soutien indéfectible tout au long de mon parcours.\\
Ce travail est le fruit de votre confiance en moi.

À ma famille,\\
dont l'encouragement et la chaleur ont été ma force\\
dans les moments les plus difficiles.

À mes amis et à tous ceux qui ont contribué,\\
de près ou de loin, à l'accomplissement de ce projet.

\vspace{1em}
\begin{flushright}
\textbf{Marwa Abada}
\end{flushright}
\end{minipage}
\end{center}

\vfill\null
\newpage

% ─────────────────────────────────────────────────────────────────────────────
% REMERCIEMENTS
% ─────────────────────────────────────────────────────────────────────────────
\thispagestyle{empty}
\begin{center}
  {\Large\bfseries Remerciements}\\[0.6cm]
\end{center}
\vspace{0.8cm}

Au terme de ce travail, il m'est particulièrement agréable d'exprimer ma
gratitude et mes sincères remerciements à toutes les personnes qui ont
contribué, de près ou de loin, à la réalisation de ce projet.

Je tiens tout d'abord à adresser mes plus vifs remerciements à
\textbf{Mme Thouraya Gouasmi}, mon encadrante, pour son accompagnement,
sa disponibilité, ses précieux conseils ainsi que son soutien tout au long
de cette expérience. Ses orientations techniques et pédagogiques ont été
d'une grande importance dans la réussite de ce projet.

J'exprime également ma profonde reconnaissance à l'ensemble des enseignants
et du personnel de \textbf{L'Institut Supérieur des Sciences Appliquées et de Technologie de Gafsa},
pour la qualité de la formation dispensée ainsi que pour les connaissances
acquises durant mon parcours universitaire.

Je remercie chaleureusement l'administration et toute l'équipe de
l'institut pour l'accueil, l'encadrement et les conditions favorables
offertes durant la période de stage.

Mes sincères remerciements vont également à ma famille et à mes proches
pour leur soutien moral, leurs encouragements constants et leur confiance
tout au long de mon parcours académique.

Enfin, je remercie toutes les personnes qui ont contribué, directement ou
indirectement, à l'aboutissement de ce travail.

\vspace{1cm}
\begin{flushright}
\itshape Marwa Abada\\
Gafsa, 2026
\end{flushright}

\newpage

% ─────────────────────────────────────────────────────────────────────────────
% RÉSUMÉ
% ─────────────────────────────────────────────────────────────────────────────
\chapter*{Résumé}
\addcontentsline{toc}{chapter}{Résumé}

Ce mémoire présente un pipeline Big Data end-to-end pour la prédiction des maladies cardiovasculaires (MCV) à partir du \textit{Behavioral Risk Factor Surveillance System} (BRFSS) 2020--2024, comptant plus de 2,17 millions d'enregistrements épidémiologiques américains.

Le pipeline traite 2\,176\,776 enregistrements BRFSS bruts avec Apache Spark (PySpark)~: déduplication de 416\,535 lignes, binarisation des codes BRFSS non standards, et imputation des valeurs manquantes. Un feature engineering multi-méthodes (corrélations de Pearson et Spearman, importance Gini, test du $\chi^2$) sélectionne 15 variables parmi lesquelles deux features dérivées composites, RiskScore et HealthIndex. Deux modèles sont entraînés sur un partitionnement stratifié 70/15/15~: un réseau de neurones profond à architecture 256-128-64 (TensorFlow~2.19) et XGBoost~3.2 avec \textit{RandomizedSearchCV} sur 20 combinaisons de paramètres.

Les deux modèles convergent vers une AUC-ROC de 0,819 (DNN~: 0,8191, XGBoost~: 0,8195). Le DNN, avec un seuil de décision optimisé à 0,6713 par courbe Précision-Rappel, obtient une accuracy de 82,25~\% et un recall de 56,3~\% sur la classe positive. XGBoost, avec 400 arbres et un \textit{scale\_pos\_weight} de 8,69, atteint un recall de 79,9~\%. Le benchmark sur cinq fractions croissantes du jeu d'entraînement montre qu'XGBoost standalone est 36 à 97 fois plus rapide que Spark~GBT en mode local~; l'apport de Spark dans ce pipeline est donc concentré sur l'étape de preprocessing distribué.

Un tableau de bord interactif Streamlit a été développé pour visualiser l'ensemble des résultats.

\vspace{0.4cm}
\textbf{Mots-clés :} Maladies cardiovasculaires, Big Data, Apache Spark, Deep Learning, XGBoost, BRFSS, Classification binaire, Déséquilibre de classes.

\newpage

% ─────────────────────────────────────────────────────────────────────────────
% ABSTRACT
% ─────────────────────────────────────────────────────────────────────────────
\chapter*{Abstract}
\addcontentsline{toc}{chapter}{Abstract}

This thesis presents an end-to-end Big Data pipeline for cardiovascular disease (CVD) prediction using the Behavioral Risk Factor Surveillance System (BRFSS) 2020--2024 dataset, comprising over 2.17 million epidemiological records.

The work is organized around five main axes: (1) data ingestion and cleaning with Apache Spark, processing 2.17 million rows and removing 416,535 duplicates; (2) feature engineering including Pearson/Spearman correlation analysis, Gini importance, chi-squared testing, and creation of derived features (RiskScore, HealthIndex); (3) training a Deep Neural Network (DNN) with a 256-128-64 architecture and class-weight-based imbalance handling; (4) training and optimizing XGBoost via RandomizedSearchCV; and (5) a performance benchmark between XGBoost standalone and Spark GBTClassifier across increasing dataset fractions.

Results show that both models achieve comparable AUC-ROC: 0.8191 for the DNN and 0.8195 for XGBoost. Using a threshold optimized at 0.6713 via the Precision-Recall curve, the DNN achieves 82.25\% accuracy and 56.3\% recall on the positive class. XGBoost, with 400 trees and a \textit{scale\_pos\_weight} of 8.69, reaches 79.9\% recall. The benchmark demonstrates that XGBoost standalone is 36 to 97 times faster than Spark GBT in local single-machine mode, highlighting that Spark's contribution lies in distributed preprocessing rather than model training. An interactive Streamlit dashboard was developed to visualize all results.

\vspace{0.4cm}
\textbf{Keywords:} Cardiovascular diseases, Big Data, Apache Spark, Deep Learning, XGBoost, BRFSS, Binary classification, Class imbalance.

\newpage
